All approximation in this chapter will use in some way the Lemma of Laplace, or the Laplace Method for integrals of the form
\begin{equation}
	\Phi(\lambda) \deff \int_{a}^{b} g(x) \ee^{-\lambda \phi(x)} \dd x
	\label{Equation: Laplace}
\end{equation}
We first state a non-precise version of this theorem and discuss some ideas surrounding it.
\begin{thm}
	Let $\Phi$ be as in \eqref{Equation: Laplace}, with $g(x)$ a continuous function and $f(x)$ assuming a global maximum at $x^* \in (a,b)$ with $f'(x^*) = 0$ and $f''(x^*) > 0$. Then, under some regularity conditions, the function $\Phi$ satisfies an estimate of the form
	\begin{equation}
		\Phi(\lambda) = \ee^{-\lambda \phi(x^*)} \frac{\sqrt{2\pi} g(x^*)}{\sqrt{|\phi''(x^*)|}\lambda^{1/2}} \left( 1 + \boh(1)\right)
	\end{equation}
	as $\lambda \rightarrow \infty$.
	\label{Prop: Non-formal Laplace's Method}
\end{thm}

The idea of the proof is quite more revealing than the proper result. We abstain for now from giving a proper formal proof and instead give a general sketch that should be enough to illustrate the important point for this section.

\begin{proof}
	(Sketch of Proof) The main idea is to approximate $\phi$ around its global maximum by Taylor's theorem 
	\begin{equation}
		\phi(x) = \phi(x^*) - \frac{1}{2} |\phi''(x^*)|(x-x^*)^2 + \epsilon_2(x-x^*)
	\end{equation}
	and then perform a change of variables $\frac{s}{\sqrt{\lambda}} = x - x^*$, so that 
	\begin{equation}
		\int_{x^*-\frac{c}{\sqrt{\lambda}}}^{x^*+\frac{c}{\sqrt{\lambda}}} \ee^{-\lambda \phi(x)} \dd(x) = \ee^{-\lambda \phi(x^*)} \int_{-c}^c \ee^{\frac{1}{2} |\phi''(x^*)|t^2 - \lambda \epsilon_2\left(\frac{t}{\sqrt{\lambda}}\right)} \frac{\dd t}{\sqrt{\lambda}}.
	\end{equation}
	Now, using the technical assumption of regularity on the function $\phi$, one can argue that the error term $\epsilon_2$ goes to zero sufficiently fast, leaving a Gaussian Integral. More that that, one need to argue that the integral outside this small neighborhood of $x^*$ is exponentially smaller then $\ee^{\lambda \phi(x^*)}$ and use continuity of $g$ to include it in the argument.
\end{proof}

Some important points to highlight:
\begin{enumerate}
	\item More precise error estimates can be obtained if one uses higher order terms of the Taylor Series expansion of $\phi$.
	\item This argument and result could also be used on curves on the plane. This is clear from the fact that every (smooth) curve is locally a straight line. This will lead us to the Method of \textit{Steepest Descent}.
	\item For maximizer $x^*$ which have null second derivative (a double critical point), the Laplace Approximation can be used expanding to the third order terms of the Taylor Series $\phi(x) = \phi(x^*) = \frac{1}{6}\phi'''(x^*)(x-x^*)^3 + \epsilon_3(x-x^*)$ and making the change of variables $\frac{t}{\sqrt[3]{n}} = x - x^*$ instead. This is precisely one of the cases in the expansion for the Hermite Polynomial. 
\end{enumerate}

These next subsections are in the spirit of proving a bit more formally the expansion given by the Laplace Method. 

\subsection{Watson's Lemma}
% Section Watson Lemma
\input{./Text/Subsections/WatsonLemma.tex}

\subsection{Laplace's Lemma}
% Section Laplace Lemma
\input{./Text/Subsections/LaplaceLemma.tex}